\section*{Resumo}

A biometria é uma tecnologia de identificação que se baseia em características físicas e comportamentais únicas de cada indivíduo, como impressões digitais, íris e voz. Esses dados, quando capturados por sensores, podem ser convertidos em representações numéricas e processados por algoritmos de reconhecimento de padrões. Por serem intrínsecos e exclusivos, os dados biométricos oferecem um meio de autenticação mais seguro e confiável do que métodos tradicionais, como senhas. Este projeto foca na vertente física, especificamente na identificação por impressões digitais, que utiliza os padrões formados por elevações e depressões na pele das pontas dos dedos. Sistemas biométricos funcionam comparando características extraídas desses padrões com informações previamente armazenadas em bancos de dados. Um dos principais desafios nesse processo está na variação natural entre diferentes capturas de uma mesma digital, devido a fatores como iluminação, rotação ou pressão. Assim, torna-se essencial a extração de características e formas de representação que permitam diferenciar com precisão diferentes digitais, ao mesmo tempo em que mantenham consistência para uma mesma digital. O objetivo deste projeto é avaliar diferentes técnicas de extração de características para o reconhecimento de usuários pela impressão digital.
